\chapter{Diagrama de Conexões e Mapeamento de Pinos do ESP32}
\label{anexo-pinagem}

O Quadro~\ref{qua:pinagem} apresenta o mapeamento de pinagem física planejado para o microcontrolador ESP32 (placa NodeMCU-32S) para conexão dos sensores, circuitos de proteção e atuador vibrotátil.

\begin{quadro}[htb]
\centering
\caption{Mapeamento dos pinos de entrada e saída (GPIO) previsto para o protótipo}
\label{qua:pinagem}
\small
\begin{tabular}{|l|l|l|p{7.5cm}|}
\hline
\textbf{Pino ESP32} & \textbf{Direção} & \textbf{Nível} & \textbf{Função e Dispositivo Conectado} \\
\hline
GPIO 5 & Saída & 3,3~V & Disparo do sinal de ultrassom (\textit{Trigger}) do HC-SR04. \\
\hline
GPIO 18 & Entrada & 3,3~V & Recepção do eco (\textit{Echo}) via divisor resistivo ($1\text{ k}\Omega / 2\text{ k}\Omega$). \\
\hline
GPIO 19 & Saída (PWM) & 3,3~V & Controle de modulação háptica na base do transistor NPN 2N2222. \\
\hline
VIN / 5V & Alimentação & 5,0~V & Barramento de alimentação do sensor HC-SR04 e do coletor do motor. \\
\hline
GND & Referência & 0~V & Referência elétrica comum a todos os módulos e atuadores. \\
\hline
\end{tabular}
\normalsize
\fonte{Elaborado pelos autores (2026).}
\end{quadro}

\chapter{Estrutura Preliminar do Firmware Embarcado no ESP32}
\label{anexo-firmware}

Apresenta-se a seguir a estrutura preliminar de código em C++ planejada para implementação no ESP32 para a leitura do sensor ultrassônico, aplicação de filtro de média móvel e controle da resposta háptica:

\begin{verbatim}
// Protótipo de Bracelete Inteligente para Pessoas com Baixa Visão
// IFPR Campus Pinhais - Práticas de Extensão III (2026)
// Autores: Matheus Vitor Lourenço Schionato e Vairtles Nehisen M. Liel

const int PIN_TRIG = 5;
const int PIN_ECHO = 18;
const int PIN_MOTOR = 19;

const int JANELA_FILTRO = 5;
float historico[JANELA_FILTRO] = {0};
int indiceAmostra = 0;

void setup() {
  Serial.begin(115200);
  pinMode(PIN_TRIG, OUTPUT);
  pinMode(PIN_ECHO, INPUT);
  pinMode(PIN_MOTOR, OUTPUT);
  digitalWrite(PIN_TRIG, LOW);
  digitalWrite(PIN_MOTOR, LOW);
}

float medirDistancia() {
  digitalWrite(PIN_TRIG, LOW);
  delayMicroseconds(2);
  digitalWrite(PIN_TRIG, HIGH);
  delayMicroseconds(10);
  digitalWrite(PIN_TRIG, LOW);

  long duracao = pulseIn(PIN_ECHO, HIGH, 30000); // timeout 30ms (~5m)
  if (duracao == 0) return 400.0; // Sem eco = livre

  float distancia = (duracao * 0.0343) / 2.0;
  return distancia;
}

float aplicarFiltro(float novaLeitura) {
  historico[indiceAmostra] = novaLeitura;
  indiceAmostra = (indiceAmostra + 1) % JANELA_FILTRO;
  float soma = 0;
  for (int i = 0; i < JANELA_FILTRO; i++) {
    soma += historico[i];
  }
  return soma / JANELA_FILTRO;
}

void loop() {
  float dBruta = medirDistancia();
  float dFiltrada = aplicarFiltro(dBruta);

  Serial.printf("Distancia: %.1f cm\n", dFiltrada);

  if (dFiltrada <= 60.0) {
    // Zona Crítica: Vibração contínua imediata
    digitalWrite(PIN_MOTOR, HIGH);
  } else if (dFiltrada <= 120.0) {
    // Zona de Atenção: Pulso intermitente de alerta
    digitalWrite(PIN_MOTOR, HIGH);
    delay(150);
    digitalWrite(PIN_MOTOR, LOW);
    delay(350);
  } else {
    // Zona Livre: Sem vibração (repouso)
    digitalWrite(PIN_MOTOR, LOW);
  }
  delay(50); // Taxa de atualização ~20 Hz
}
\end{verbatim}